\section*{Bab \ref{ch:g11:func} --- Fungsi dan Variasinya}

\begin{solution}{exo:g11:func:1}
$f$: menuntut $3 - x \geq 0$, \omterm{def:g11:func:function}{daerah asalnya} $\intoc{-\infty}{3}$.

$g$: menuntut $x^2 - 4 \neq 0$, yaitu $x \neq \pm 2$: \omterm{def:g11:func:function}{daerah asalnya}
$\R \setminus \{-2, 2\}$.

$h$: menuntut $x \geq 0$ dan $x \neq 5$: \omterm{def:g11:func:function}{daerah asalnya}
$\intco{0}{5} \cup \intoo{5}{+\infty}$.
\end{solution}

\begin{solution}{exo:g11:func:2}
$\sqrt{11} < \sqrt{13}$ sebab $\sqrt{}$ bersifat \omterm{def:g11:func:monotone}{naik} dan $11 < 13$.

$(-2.1)^2 > (-2.05)^2$ sebab $x^2$ turun pada
$\intoc{-\infty}{0}$ dan $-2.1 < -2.05$.

$\frac{1}{0.9} > \frac{1}{1.1}$ sebab $\frac1x$ turun pada
$\intoo{0}{+\infty}$ dan $0.9 < 1.1$.

$(-1.01)^3 < (-0.99)^3$ sebab $x^3$ \omterm{def:g11:func:monotone}{naik} dan $-1.01 < -0.99$.
\end{solution}

\begin{solution}{exo:g11:func:3}
$f(-x) = 3(-x)^4 - (-x)^2 = 3x^4 - x^2 = f(x)$: \omterm{def:g11:func:parity}{genap}.

$g(-x) = -x^3 - x = -g(x)$: \omterm{def:g11:func:parity}{ganjil}.

$h(1) = 2$ dan $h(-1) = 1 - 1 = 0$, jadi $h(-1)$ bukan $h(1)$ dan bukan
pula $-h(1)$: bukan \omterm{def:g11:func:parity}{genap} maupun \omterm{def:g11:func:parity}{ganjil}.

$k(-x) = \frac{-x}{x^2+1} = -k(x)$: \omterm{def:g11:func:parity}{ganjil}.
\end{solution}

\begin{solution}{exo:g11:func:4}
Bacalah \omterm{def:g10:functions:table}{tabel variasinya}. Pada $\intcc{-3}{1}$, $f$ mendaki dari $-2$ ke $4$
dan melewati nilai $1$ tepat sekali; pada $\intcc{1}{5}$ \omterm{def:g11:func:function}{fungsi} itu menurun
dari $4$ ke $0$ dan melewati $1$ sekali lagi: jadi \omterm{def:g10:algebra:equation}{persamaan} $f(x) = 1$
mempunyai $2$ penyelesaian. \omterm{def:g10:functions:extrema}{Maksimum} $f$ adalah $f(1) = 4 < 5$, sehingga
$f(x) = 5$ tidak mempunyai penyelesaian.
\end{solution}

\begin{solution}{exo:g11:func:5}
Misalkan $2 \leq u < v$. Maka
\[
f(v) - f(u) = v^2 - 4v - u^2 + 4u = (v-u)(v+u) - 4(v-u)
= (v - u)(u + v - 4).
\]
Kedua faktornya positif: $v - u > 0$, dan $u + v > 2 + 2 = 4$ sebab
$u \geq 2$ dan $v > 2$. Jadi $f(v) > f(u)$: $f$ tegas \omterm{def:g11:func:monotone}{naik}
pada $\intco{2}{+\infty}$.
\end{solution}

\begin{solution}{exo:g11:func:6}
$f = \sqrt{u}$ dengan $u(x) = 5 - x$ yang turun dan taknegatif pada
$\intoc{-\infty}{5}$: jadi $f$ turun di sana
(\cref{prop:g11:func:transforms}).

$g = \frac1u$ dengan $u(x) = x^2 + 1$ yang \omterm{def:g11:func:monotone}{naik} dan positif pada
$\intco{0}{+\infty}$: jadi $g$ turun di sana.

$h = 3 - 2\sqrt x$: $\sqrt x$ \omterm{def:g11:func:monotone}{naik}, mengalikannya dengan $-2 < 0$
membalik arahnya, dan menambahkan $3$ tidak mengubah apa pun: jadi $h$ turun
pada $\intco{0}{+\infty}$.
\end{solution}

\begin{solution}{exo:g11:func:7}
\emph{1.} $f$ sudah dalam \omterm{thm:g11:quad:canonical}{bentuk kanonik}: turun pada $\intoc{-\infty}{2}$,
\omterm{def:g11:func:monotone}{naik} pada $\intco{2}{+\infty}$, dengan \omterm{def:g10:functions:extrema}{minimum} $f(2) = -3$.

\emph{2.} $g = \frac{1}{u}$ dengan $u(x) = (x-2)^2 + 1 = f(x) + 4$.
Pergeseran sebesar $4$ mempertahankan variasi $f$, dan $u \geq 1 > 0$.
Mengambil kebalikannya membalik variasinya: $g$ \omterm{def:g11:func:monotone}{naik} pada
$\intoc{-\infty}{2}$ dan turun pada $\intco{2}{+\infty}$, dengan \omterm{def:g10:functions:extrema}{maksimum}
$g(2) = \frac{1}{1} = 1$.
\end{solution}

\begin{solution}{exo:g11:func:8}
Sebarang kurva berbentuk bukit yang melalui ketiga titik itu bisa dipakai. Lalu:
$f + 2$ melalui $(0,3)$, $(2,5)$, $(4,3)$, dengan variasi yang sama
(pergeseran tegak); $-f$ melalui $(0,-1)$, $(2,-3)$, $(4,-1)$ dengan
variasi yang \emph{terbalik} (puncaknya menjadi lembah); $2f$ melalui
$(0,2)$, $(2,6)$, $(4,2)$, dengan variasi yang sama tetapi teregang
secara tegak.
\end{solution}

\begin{solution}{exo:g11:func:9}
$2 - \frac{1}{x+1} = \frac{2(x+1) - 1}{x+1} = \frac{2x+1}{x+1} = f(x)$.
Pada $\intoo{-1}{+\infty}$, $u(x) = x + 1$ \omterm{def:g11:func:monotone}{naik} dan positif, sehingga
$\frac1u$ turun, $-\frac1u$ \omterm{def:g11:func:monotone}{naik}, dan menambahkan $2$ tidak mengubah
apa pun: jadi $f$ tegas \omterm{def:g11:func:monotone}{naik} pada $\intoo{-1}{+\infty}$.
\end{solution}

\begin{solution}{exo:g11:func:10}
\emph{1. Benar.} Jika $u < v$, maka
$(f+g)(v) - (f+g)(u) = \bigl(f(v)-f(u)\bigr) + \bigl(g(v)-g(u)\bigr)$
merupakan jumlah dua bilangan taknegatif.

\emph{2. Salah.} $f(x) = g(x) = x$ \omterm{def:g11:func:monotone}{naik} pada $\R$, tetapi hasil
kalinya $x^2$ tidak \omterm{def:g11:func:monotone}{naik} pada $\R$ (\omterm{def:g11:func:function}{fungsi} itu turun pada
$\intoc{-\infty}{0}$). (Pernyataannya menjadi benar bila kedua \omterm{def:g11:func:function}{fungsinya}
juga taknegatif.)

\emph{3. Salah.} Contoh penyangkal yang sama: $f(x) = x$ \omterm{def:g11:func:monotone}{naik} pada $\R$
tetapi $f^2(x) = x^2$ tidak.
\end{solution}

\begin{solution}{exo:g11:func:11}
\emph{1.} Untuk $0 < u < v$:
\[
f(v) - f(u) = (v - u) + \frac1v - \frac1u
= (v - u) + \frac{u - v}{uv}
= (v - u)\left(1 - \frac{1}{uv}\right)
= (v - u)\,\frac{uv - 1}{uv}.
\]

\emph{2.} Selalu berlaku $v - u > 0$ dan $uv > 0$. Jika $0 < u < v \leq 1$,
maka $uv < 1$, jadi $uv - 1 < 0$ dan $f(v) < f(u)$: $f$ tegas turun
pada $\intoc{0}{1}$. Jika $1 \leq u < v$, maka $uv > 1$ dan $f(v) > f(u)$:
$f$ tegas \omterm{def:g11:func:monotone}{naik} pada $\intco{1}{+\infty}$.

\emph{3.} $f$ turun sebelum $1$ dan \omterm{def:g11:func:monotone}{naik} sesudahnya, sehingga
$f(x) \geq f(1) = 2$ untuk semua $x > 0$, dengan kesamaan tepat di $x = 1$.
(Ketaksamaan ini muncul kembali dalam banyak rupa; bandingkan
\cref{ex:g11:quad:canonical} dan ketaksamaan garis singgung pada
\cref{ch:g12:deriv}.)
\end{solution}

\begin{solution}{pb:g11:func:1}
\textbf{1.} $x^4 - 3x^2$: \omterm{def:g11:func:parity}{genap} ($(-x)^4 - 3(-x)^2 = x^4 -
3x^2$). $x^3 - x$: \omterm{def:g11:func:parity}{ganjil}. $x^2 + x$: bukan keduanya ($f(-1) = 0 \neq
f(1) = 2$ dan $\neq -f(1)$). $\abs x$: \omterm{def:g11:func:parity}{genap}. $\frac1x$: \omterm{def:g11:func:parity}{ganjil}.

\textbf{2.} \omterm{def:g11:func:parity}{Genap}: kurvanya setangkup terhadap sumbu $y$
(cermin). \omterm{def:g11:func:parity}{Ganjil}: setangkup terhadap titik asal (setengah putaran).

\textbf{3.} Jika $f$ dan $g$ \omterm{def:g11:func:parity}{ganjil}:
$(fg)(-x) = f(-x)g(-x) = (-f(x))(-g(x)) = (fg)(x)$: jadi \omterm{def:g11:func:parity}{genap}.
Tabelnya: \omterm{def:g11:func:parity}{genap} $\times$ \omterm{def:g11:func:parity}{genap} $=$ \omterm{def:g11:func:parity}{genap}; \omterm{def:g11:func:parity}{genap} $\times$ \omterm{def:g11:func:parity}{ganjil} $=$ \omterm{def:g11:func:parity}{ganjil};
\omterm{def:g11:func:parity}{ganjil} $\times$ \omterm{def:g11:func:parity}{ganjil} $=$ \omterm{def:g11:func:parity}{genap} --- yaitu aturan tanda dengan paritas.

\textbf{4.} \omterm{def:g11:func:parity}{Ganjil} di $0$: $f(0) = -f(0)$, jadi $f(0) = 0$. \omterm{def:g11:func:parity}{Genap}:
$f(-3) = f(3) = 7$. Sekaligus \omterm{def:g11:func:parity}{genap} dan \omterm{def:g11:func:parity}{ganjil}:
$f(x) = f(-x) = -f(x)$ untuk semua $x$, jadi $2f(x) = 0$: $f$ adalah
\omterm{def:g11:func:function}{fungsi} nol.

\textbf{5.} Paruh kanannya ($x \geq 0$) sudah menentukan sebuah \omterm{def:g11:func:parity}{fungsi
genap} sepenuhnya: paruh kirinya adalah cerminannya. Bukan \omterm{def:g11:func:parity}{genap}
maupun \omterm{def:g11:func:parity}{ganjil}: $x^2 + x$ (pertanyaan 1) --- kedua kesimetriannya gagal di
$x = 1$.

\textbf{6.} $E(x) = \frac{(x^3 + x^2 + 1) + (-x^3 + x^2 +
1)}{2} = x^2 + 1$ (\omterm{def:g11:func:parity}{genap}), $O(x) = \frac{(x^3 + x^2 + 1) -
(-x^3 + x^2 + 1)}{2} = x^3$ (\omterm{def:g11:func:parity}{ganjil}); jumlahnya adalah $f$.

\textbf{7.} $E(-x) = \frac{f(-x) + f(x)}{2} = E(x)$: \omterm{def:g11:func:parity}{genap}.
$O(-x) = \frac{f(-x) - f(x)}{2} = -O(x)$: \omterm{def:g11:func:parity}{ganjil}. Lalu
$E(x) + O(x) = f(x)$ lewat penjumlahan langsung: jadi setiap \omterm{def:g11:func:function}{fungsi}
terbelah.

\textbf{8.} Dari $E_1 + O_1 = E_2 + O_2$:
$E_1 - E_2 = O_2 - O_1$. Ruas kirinya \omterm{def:g11:func:parity}{genap}, ruas kanannya
\omterm{def:g11:func:parity}{ganjil} --- jadi \omterm{def:g11:func:function}{fungsi} itu sekaligus keduanya, sehingga nol (pertanyaan 4):
$E_1 = E_2$ dan $O_1 = O_2$. Hanya ada satu pembelahan.

\textbf{9.} $E(x) = \frac12\left(\frac{1}{1 - x} + \frac{1}{1 +
x}\right) = \frac12 \cdot \frac{(1 + x) + (1 - x)}{1 - x^2}
= \frac{1}{1 - x^2}$, dan
$O(x) = \frac12 \cdot \frac{(1 + x) - (1 - x)}{1 - x^2}
= \frac{x}{1 - x^2}$. Periksa: $E$ \omterm{def:g11:func:parity}{genap}, $O$ \omterm{def:g11:func:parity}{ganjil}, dan
jumlahnya $\frac{1 + x}{1 - x^2} =
\frac{1 + x}{(1 - x)(1 + x)} = \frac{1}{1 - x} = f(x)$.

\textbf{10.} Karena setiap potongnya menuruti sebuah kesimetrian, separuh
pekerjaan (dan separuh datanya) sudah cukup untuk masing-masing: selesaikan
di satu sisi, lalu cerminkan ke sisi lain --- dan banyak operasi
(pengintegralan, deret, analisis isyarat) memperlakukan potongan yang
setangkup jauh lebih sederhana daripada \omterm{def:g11:func:function}{fungsi} mentah.

\textbf{11.} $(x-3)^2$: tergeser $3$ ke kanan. $x^2 + 2$: tergeser
$2$ ke atas. $(x-3)^2 + 2$: keduanya. $-x^2$: terbalik ke bawah.
$2x^2$: teregang tegak sebesar $2$ (cekungannya lebih sempit).

\textbf{12.} Untuk $x \geq 2$: $(x - 2) + (x + 2) = 2x$. Untuk
$-2 \leq x \leq 2$: $(2 - x) + (x + 2) = 4$. Untuk $x \leq -2$:
$-2x$. Kurvanya menurun dengan \omterm{def:g10:reffunc:affine}{gradien} $-2$, mendatar pada tinggi
$4$ antara $-2$ dan $2$, lalu mendaki dengan \omterm{def:g10:reffunc:affine}{gradien} $2$: sebuah lembah
berlantai rata.

\textbf{13.} Pergeseran tegak: $f(x) + 2$ tetap \omterm{def:g11:func:parity}{genap}
($f(-x) + 2 = f(x) + 2$). Peregangan tegak: $2f(x)$ tetap
\omterm{def:g11:func:parity}{genap}. Pergeseran mendatar: $(x - 3)^2$ tidak \omterm{def:g11:func:parity}{genap}
($f(-1) = 16 \neq 4 = f(1)$): bergeser ke samping merusak
cerminnya.

\textbf{14.} Pada lantai ratanya nilainya $4 \neq 6$. Untuk
$x \geq 2$: $2x = 6$, $x = 3$. Untuk $x \leq -2$: $-2x = 6$,
$x = -3$. Penyelesaiannya: $\pm 3$.

\textbf{15.} $f(x) = a(x - 2)^2 - 3$ dengan $f(0) = 4a - 3 = 1$:
jadi $a = 1$, sehingga $f(x) = (x - 2)^2 - 3$.

\textbf{16.} $g(v) - g(u) = (v - u) + 4\,\frac{u - v}{uv}
= (v - u)\,\frac{uv - 4}{uv}$: bernilai negatif untuk $u < v$ dengan
$uv < 4$, dan positif dengan $uv > 4$: jadi turun pada
$\intoc{0}{2}$, \omterm{def:g11:func:monotone}{naik} pada $\intco{2}{+\infty}$; \omterm{def:g10:functions:extrema}{minimumnya} di
$x = 2$, bernilai $g(2) = 4$.

\textbf{17.} Sisinya $x$ dan $\frac{16}{x}$:
$P(x) = 2\left(x + \frac{16}{x}\right)$, yang \omterm{def:g10:functions:extrema}{minimum} (dengan mesin
pertanyaan 16 dan $4 \to 16$, titik baliknya $\sqrt{16} = 4$)
di $x = 4$: yaitu \emph{persegi} $4 \times 4$, kelilingnya $16$ ---
jawaban Dido, kini terbukti untuk setiap panjang sisi real.

\textbf{18.} Pada $\intco{0}{+\infty}$: $x \mapsto x^2 + 1$
\omterm{def:g11:func:monotone}{naik}, dan $\sqrt{\phantom{x}}$ \omterm{def:g11:func:monotone}{naik}, sehingga
komposisinya \omterm{def:g11:func:monotone}{naik}. Lalu $h(-x) = \sqrt{x^2 + 1} = h(x)$: $h$
\omterm{def:g11:func:parity}{genap}, sehingga menurut kesimetrian cerminnya \omterm{def:g11:func:function}{fungsi} itu turun pada
$\intoc{-\infty}{0}$. \omterm{def:g10:functions:extrema}{Minimumnya} $h(0) = 1$.

\textbf{19.} Kuadrat jarak ke $(3, 0)$ dari titik
$(x, \sqrt x)$ adalah
$d^2(x) = (x - 3)^2 + x = x^2 - 5x + 9$, yang \omterm{def:g10:functions:extrema}{minimum} di
$x = \frac52$: jadi titik terdekatnya
$\left(\frac52, \sqrt{\frac52}\right)$, dengan jarak
$\sqrt{\frac{11}{4}} = \frac{\sqrt{11}}{2}$. Meminimumkan $d^2$
sah karena pengkuadratan tegas \omterm{def:g11:func:monotone}{naik} pada
bilangan positif: $d$ dan $d^2$ mencapai dasar di tempat yang sama.

\textbf{20.} Bahan mentahnya: kurva acuan, yang dihafal di luar kepala.
Mesinnya: pergeseran, pembalikan dan peregangan yang merangkai \omterm{def:g10:functions:graph}{grafik} baru
dari yang lama. Kendali mutunya: paritas dan pembelahan genap--ganjil,
yang memangkas separuh setiap telaah. Bangku ukurnya: pemfaktoran
$f(v) - f(u)$, yang membaca variasi dan \omterm{def:g10:functions:extrema}{minimum} secara langsung. Perkakas
ampuh bab berikutnya: \emph{turunan}, yang membaca semuanya sekilas.

\end{solution}
